\documentclass[11pt,a4paper]{article}
\usepackage[utf8]{inputenc}
\usepackage[T1]{fontenc}
\usepackage[spanish,es-minimal,es-noshorthands]{babel}
\usepackage{lmodern}
\usepackage[sfdefault,lining]{FiraSans}
\usepackage[a4paper,top=2.3cm,bottom=2.2cm,left=1.8cm,right=1.8cm,headheight=1.3cm,headsep=0.4cm,footskip=1cm]{geometry}
\usepackage{amsmath,amssymb}
\usepackage{icomma}
\usepackage{graphicx}
\usepackage[table]{xcolor}
\usepackage{tikz}
\usepackage{fontawesome5}
\usepackage{ccicons}
\usepackage{fancyhdr}
\usepackage{lastpage}
\usepackage{tasks}
\usepackage{needspace}
\usepackage[most]{tcolorbox}
\usepackage{array}
\usepackage[hidelinks]{hyperref}
\AtBeginDocument{\hypersetup{pdftitle={Ficha de ejercicios de la unidad 4: Fracciones. 1º ESO},pdfsubject={Matemáticas 1º ESO},pdfkeywords={fracciones, fracción de un número, porcentajes, fracciones equivalentes, común denominador, comparar fracciones, operaciones con fracciones, problemas}}}

\definecolor{navy}{HTML}{1B2748}
\definecolor{accent}{HTML}{E8603A}
\definecolor{teal}{HTML}{1F9E90}
\definecolor{soft}{HTML}{EEF1F8}
\definecolor{muted}{HTML}{5D6685}

\setlength{\parindent}{0pt}
\setlength{\parskip}{3pt}
\settasks{label=\alph*), label-format=\bfseries\color{navy}, label-width=1.3em, item-indent=1.8em, column-sep=0.6em, after-item-skip=5pt, before-skip=3pt, after-skip=4pt}

% ---------- capçalera i peu ----------
\newcommand{\llicencia}{\href{https://creativecommons.org/licenses/by-nc-sa/4.0/deed.es}{\ccbyncsa}\ \ Licencia Creative Commons Atribución-NoComercial-CompartirIgual 4.0 Internacional (CC BY-NC-SA 4.0)}
\pagestyle{fancy}
\fancyhf{}
\fancyhead[L]{\raisebox{-0.35\height}{\includegraphics[height=0.95cm]{logo-gva.png}}}
\fancyhead[C]{\small\color{muted}Ficha de ejercicios · Unidad 4: Fracciones · 1\textsuperscript{o} ESO}
\fancyhead[R]{\raisebox{-0.35\height}{\includegraphics[height=0.95cm]{logo-vermellar.png}}}
\fancyfoot[L]{\scriptsize\color{muted}\llicencia}
\fancyfoot[R]{\small\color{muted}\thepage}
\renewcommand{\headrulewidth}{0.4pt}
\renewcommand{\headrule}{\color{muted!40}\hrule height \headrulewidth}
\fancypagestyle{primera}{\fancyhf{}\renewcommand{\headrulewidth}{0pt}
  \fancyfoot[L]{\scriptsize\color{muted}\llicencia}
  \fancyfoot[R]{\small\color{muted}\thepage}}

% ---------- nivells i exercicis ----------
\newcommand{\estrelles}[1]{{\color{accent}\ifcase#1\or\faStar\faStar[regular]\faStar[regular]\or\faStar\faStar\faStar[regular]\or\faStar\faStar\faStar\fi}}
\newcounter{exer}
\newcommand{\ex}[1]{\par\needspace{4\baselineskip}\medskip\stepcounter{exer}\noindent{\color{navy}\bfseries\theexer.}\ \ {\small\estrelles{#1}}\ \ }

\newcommand{\seccio}[2]{\par\needspace{10\baselineskip}\bigskip
  \begin{tcolorbox}[colback=navy,colframe=navy,coltext=white,arc=2mm,boxrule=0pt,left=3mm,right=3mm,top=1.2mm,bottom=1.2mm]
  {\large\bfseries #1\quad #2}
  \end{tcolorbox}\vspace{-2pt}}

\newcolumntype{C}[1]{>{\centering\arraybackslash}p{#1}}

% quadre de recordatori (només a la fitxa d'exercicis)
\newcommand{\recorda}[1]{\par\smallskip\begin{tcolorbox}[colback=accent!7,colframe=accent!55,arc=2mm,boxrule=0.5pt,left=2.5mm,right=2.5mm,top=1mm,bottom=1mm]\small{\color{accent}\faLightbulb[regular]}\ \textbf{Recuerda:}\ #1\end{tcolorbox}}
\newcommand{\res}[1]{#1}

% =====================================================================
% Capçalera FentMates de les fitxes: un color per a cada curs.
% Els colors són els de la pàgina inicial (LEVELS a assets/js/materials.js);
% el logo ve de _eines/logo/pdf/ (python3 _eines/logo/genera.py).
%
% Ús, al final del preàmbul de la unitat:
%   % =====================================================================
% Capçalera FentMates de les fitxes: un color per a cada curs.
% Els colors són els de la pàgina inicial (LEVELS a assets/js/materials.js);
% el logo ve de _eines/logo/pdf/ (python3 _eines/logo/genera.py).
%
% Ús, al final del preàmbul de la unitat:
%   % =====================================================================
% Capçalera FentMates de les fitxes: un color per a cada curs.
% Els colors són els de la pàgina inicial (LEVELS a assets/js/materials.js);
% el logo ve de _eines/logo/pdf/ (python3 _eines/logo/genera.py).
%
% Ús, al final del preàmbul de la unitat:
%   \input{../comu/fentmates.tex}
%   \fmcurs{1eso}
% i, al començament del document:
%   \capcalera{Matemàtiques}{1r ESO}{Fitxa d'exercicis · Unitat 2}{Divisibilitat}
% Opcional abans de \input: \def\fmneed{14} (espai mínim abans de cada secció).
% =====================================================================
\makeatletter
\def\fm@logodir{../../logo/pdf}
\providecommand{\fmneed}{10}
\def\fm@cursnom{}

% clau del curs, color i color del text damunt del color
\newcommand{\fm@curs}[3]{\@namedef{fm@color@#1}{#2}\@namedef{fm@text@#1}{#3}}
\fm@curs{1eso}{E6B84A}{navy}
\fm@curs{2eso}{3B82F6}{white}
\fm@curs{3eso}{FF7A50}{white}
\fm@curs{4eso-a}{A78BFA}{white}
\fm@curs{4eso-b}{6D28D9}{white}
\fm@curs{1batx-mi}{9B1C1C}{white}
\fm@curs{1batx-ccss}{EF6B6B}{white}
\fm@curs{2batx-mii}{1E6B47}{white}
\fm@curs{2batx-ccss}{2BB3A3}{white}

\newcommand{\fmcurs}[1]{%
  \@ifundefined{fm@color@#1}{\PackageError{fentmates}{Curs desconegut: #1}{}}{}%
  \edef\fm@tmp{\noexpand\definecolor{curs}{HTML}{\@nameuse{fm@color@#1}}%
               \noexpand\colorlet{curstext}{\@nameuse{fm@text@#1}}}\fm@tmp
  \colorlet{cursfosc}{curs!72!black}%
  \def\fm@logo{\fm@logodir/vertical-#1.pdf}%
  \def\fm@marca{\fm@logodir/marca-#1.pdf}}

% ---------- capçalera de la primera pàgina ----------
% #1 matèria (petit), #2 curs, #3 tipus de document i unitat, #4 títol
\newcommand{\capcalera}[4]{%
  \gdef\fm@cursnom{#2}%
  \par\noindent
  \begin{tcolorbox}[colback=curs,colframe=curs,boxrule=0pt,arc=2.2mm,width=3.5cm,height=2.15cm,
      valign=center,left=3mm,right=2mm,top=1mm,bottom=1mm,nobeforeafter,box align=center]
    \color{curstext}{\footnotesize #1}\par\vspace{1pt}{\fontsize{19}{21}\selectfont\bfseries #2\par}%
  \end{tcolorbox}\hspace{5mm}%
  \begin{minipage}[c]{\dimexpr\linewidth-3.5cm-5mm-2.9cm\relax}
    {\color{cursfosc}\small\bfseries\MakeUppercase{#3}}\par\vspace{2pt}
    {\color{navy}\fontsize{20}{23}\selectfont\bfseries #4\par}
  \end{minipage}\hfill
  \begin{minipage}[c]{2.6cm}\raggedleft\includegraphics[width=2.6cm]{\fm@logo}\end{minipage}\par
  \vspace{2pt}}

% ---------- barres de secció del color del curs ----------
\renewcommand{\seccio}[2]{\par\needspace{\fmneed\baselineskip}\bigskip
  \begin{tcolorbox}[colback=curs,colframe=curs,coltext=curstext,arc=2mm,boxrule=0pt,left=3mm,right=3mm,top=1.2mm,bottom=1.2mm]
  {\large\bfseries #1\quad #2}
  \end{tcolorbox}\vspace{-2pt}}

% ---------- estrelles de dificultat del color del curs (a Batxillerat no se'n posen) ----------
\@ifundefined{estrelles}{}{%
  \renewcommand{\estrelles}[1]{{\color{curs}\ifcase#1\or\faStar\faStar[regular]\faStar[regular]\or\faStar\faStar\faStar[regular]\or\faStar\faStar\faStar\fi}}}

% ---------- capçalera i peu de les pàgines ----------
% primera pàgina: logos institucionals al peu
\fancypagestyle{primera}{\fancyhf{}\renewcommand{\headrulewidth}{0pt}
  \fancyfoot[L]{\raisebox{-0.32\height}{\includegraphics[height=0.78cm]{logo-gva.png}}\hspace{5mm}%
    \raisebox{-0.32\height}{\includegraphics[height=0.78cm]{logo-vermellar.png}}}
  \fancyfoot[R]{\scriptsize\color{muted}\href{https://creativecommons.org/licenses/by-nc-sa/4.0/deed.ca}{\ccbyncsa}\ \ CC BY-NC-SA 4.0 · fentmates.org\qquad\small\thepage}}
% la resta: el símbol a l'esquerra i el curs a la dreta, amb el títol al centre com abans
\fancyhead[L]{\raisebox{-0.3\height}{\includegraphics[height=0.62cm]{\fm@marca}}}
\fancyhead[R]{\small\bfseries\color{cursfosc}\fm@cursnom}
\renewcommand{\headrule}{\color{curs!45}\hrule height \headrulewidth}
\makeatother

%   \fmcurs{1eso}
% i, al començament del document:
%   \capcalera{Matemàtiques}{1r ESO}{Fitxa d'exercicis · Unitat 2}{Divisibilitat}
% Opcional abans de \input: \def\fmneed{14} (espai mínim abans de cada secció).
% =====================================================================
\makeatletter
\def\fm@logodir{../../logo/pdf}
\providecommand{\fmneed}{10}
\def\fm@cursnom{}

% clau del curs, color i color del text damunt del color
\newcommand{\fm@curs}[3]{\@namedef{fm@color@#1}{#2}\@namedef{fm@text@#1}{#3}}
\fm@curs{1eso}{E6B84A}{navy}
\fm@curs{2eso}{3B82F6}{white}
\fm@curs{3eso}{FF7A50}{white}
\fm@curs{4eso-a}{A78BFA}{white}
\fm@curs{4eso-b}{6D28D9}{white}
\fm@curs{1batx-mi}{9B1C1C}{white}
\fm@curs{1batx-ccss}{EF6B6B}{white}
\fm@curs{2batx-mii}{1E6B47}{white}
\fm@curs{2batx-ccss}{2BB3A3}{white}

\newcommand{\fmcurs}[1]{%
  \@ifundefined{fm@color@#1}{\PackageError{fentmates}{Curs desconegut: #1}{}}{}%
  \edef\fm@tmp{\noexpand\definecolor{curs}{HTML}{\@nameuse{fm@color@#1}}%
               \noexpand\colorlet{curstext}{\@nameuse{fm@text@#1}}}\fm@tmp
  \colorlet{cursfosc}{curs!72!black}%
  \def\fm@logo{\fm@logodir/vertical-#1.pdf}%
  \def\fm@marca{\fm@logodir/marca-#1.pdf}}

% ---------- capçalera de la primera pàgina ----------
% #1 matèria (petit), #2 curs, #3 tipus de document i unitat, #4 títol
\newcommand{\capcalera}[4]{%
  \gdef\fm@cursnom{#2}%
  \par\noindent
  \begin{tcolorbox}[colback=curs,colframe=curs,boxrule=0pt,arc=2.2mm,width=3.5cm,height=2.15cm,
      valign=center,left=3mm,right=2mm,top=1mm,bottom=1mm,nobeforeafter,box align=center]
    \color{curstext}{\footnotesize #1}\par\vspace{1pt}{\fontsize{19}{21}\selectfont\bfseries #2\par}%
  \end{tcolorbox}\hspace{5mm}%
  \begin{minipage}[c]{\dimexpr\linewidth-3.5cm-5mm-2.9cm\relax}
    {\color{cursfosc}\small\bfseries\MakeUppercase{#3}}\par\vspace{2pt}
    {\color{navy}\fontsize{20}{23}\selectfont\bfseries #4\par}
  \end{minipage}\hfill
  \begin{minipage}[c]{2.6cm}\raggedleft\includegraphics[width=2.6cm]{\fm@logo}\end{minipage}\par
  \vspace{2pt}}

% ---------- barres de secció del color del curs ----------
\renewcommand{\seccio}[2]{\par\needspace{\fmneed\baselineskip}\bigskip
  \begin{tcolorbox}[colback=curs,colframe=curs,coltext=curstext,arc=2mm,boxrule=0pt,left=3mm,right=3mm,top=1.2mm,bottom=1.2mm]
  {\large\bfseries #1\quad #2}
  \end{tcolorbox}\vspace{-2pt}}

% ---------- estrelles de dificultat del color del curs (a Batxillerat no se'n posen) ----------
\@ifundefined{estrelles}{}{%
  \renewcommand{\estrelles}[1]{{\color{curs}\ifcase#1\or\faStar\faStar[regular]\faStar[regular]\or\faStar\faStar\faStar[regular]\or\faStar\faStar\faStar\fi}}}

% ---------- capçalera i peu de les pàgines ----------
% primera pàgina: logos institucionals al peu
\fancypagestyle{primera}{\fancyhf{}\renewcommand{\headrulewidth}{0pt}
  \fancyfoot[L]{\raisebox{-0.32\height}{\includegraphics[height=0.78cm]{logo-gva.png}}\hspace{5mm}%
    \raisebox{-0.32\height}{\includegraphics[height=0.78cm]{logo-vermellar.png}}}
  \fancyfoot[R]{\scriptsize\color{muted}\href{https://creativecommons.org/licenses/by-nc-sa/4.0/deed.ca}{\ccbyncsa}\ \ CC BY-NC-SA 4.0 · fentmates.org\qquad\small\thepage}}
% la resta: el símbol a l'esquerra i el curs a la dreta, amb el títol al centre com abans
\fancyhead[L]{\raisebox{-0.3\height}{\includegraphics[height=0.62cm]{\fm@marca}}}
\fancyhead[R]{\small\bfseries\color{cursfosc}\fm@cursnom}
\renewcommand{\headrule}{\color{curs!45}\hrule height \headrulewidth}
\makeatother

%   \fmcurs{1eso}
% i, al començament del document:
%   \capcalera{Matemàtiques}{1r ESO}{Fitxa d'exercicis · Unitat 2}{Divisibilitat}
% Opcional abans de \input: \def\fmneed{14} (espai mínim abans de cada secció).
% =====================================================================
\makeatletter
\def\fm@logodir{../../logo/pdf}
\providecommand{\fmneed}{10}
\def\fm@cursnom{}

% clau del curs, color i color del text damunt del color
\newcommand{\fm@curs}[3]{\@namedef{fm@color@#1}{#2}\@namedef{fm@text@#1}{#3}}
\fm@curs{1eso}{E6B84A}{navy}
\fm@curs{2eso}{3B82F6}{white}
\fm@curs{3eso}{FF7A50}{white}
\fm@curs{4eso-a}{A78BFA}{white}
\fm@curs{4eso-b}{6D28D9}{white}
\fm@curs{1batx-mi}{9B1C1C}{white}
\fm@curs{1batx-ccss}{EF6B6B}{white}
\fm@curs{2batx-mii}{1E6B47}{white}
\fm@curs{2batx-ccss}{2BB3A3}{white}

\newcommand{\fmcurs}[1]{%
  \@ifundefined{fm@color@#1}{\PackageError{fentmates}{Curs desconegut: #1}{}}{}%
  \edef\fm@tmp{\noexpand\definecolor{curs}{HTML}{\@nameuse{fm@color@#1}}%
               \noexpand\colorlet{curstext}{\@nameuse{fm@text@#1}}}\fm@tmp
  \colorlet{cursfosc}{curs!72!black}%
  \def\fm@logo{\fm@logodir/vertical-#1.pdf}%
  \def\fm@marca{\fm@logodir/marca-#1.pdf}}

% ---------- capçalera de la primera pàgina ----------
% #1 matèria (petit), #2 curs, #3 tipus de document i unitat, #4 títol
\newcommand{\capcalera}[4]{%
  \gdef\fm@cursnom{#2}%
  \par\noindent
  \begin{tcolorbox}[colback=curs,colframe=curs,boxrule=0pt,arc=2.2mm,width=3.5cm,height=2.15cm,
      valign=center,left=3mm,right=2mm,top=1mm,bottom=1mm,nobeforeafter,box align=center]
    \color{curstext}{\footnotesize #1}\par\vspace{1pt}{\fontsize{19}{21}\selectfont\bfseries #2\par}%
  \end{tcolorbox}\hspace{5mm}%
  \begin{minipage}[c]{\dimexpr\linewidth-3.5cm-5mm-2.9cm\relax}
    {\color{cursfosc}\small\bfseries\MakeUppercase{#3}}\par\vspace{2pt}
    {\color{navy}\fontsize{20}{23}\selectfont\bfseries #4\par}
  \end{minipage}\hfill
  \begin{minipage}[c]{2.6cm}\raggedleft\includegraphics[width=2.6cm]{\fm@logo}\end{minipage}\par
  \vspace{2pt}}

% ---------- barres de secció del color del curs ----------
\renewcommand{\seccio}[2]{\par\needspace{\fmneed\baselineskip}\bigskip
  \begin{tcolorbox}[colback=curs,colframe=curs,coltext=curstext,arc=2mm,boxrule=0pt,left=3mm,right=3mm,top=1.2mm,bottom=1.2mm]
  {\large\bfseries #1\quad #2}
  \end{tcolorbox}\vspace{-2pt}}

% ---------- estrelles de dificultat del color del curs (a Batxillerat no se'n posen) ----------
\@ifundefined{estrelles}{}{%
  \renewcommand{\estrelles}[1]{{\color{curs}\ifcase#1\or\faStar\faStar[regular]\faStar[regular]\or\faStar\faStar\faStar[regular]\or\faStar\faStar\faStar\fi}}}

% ---------- capçalera i peu de les pàgines ----------
% primera pàgina: logos institucionals al peu
\fancypagestyle{primera}{\fancyhf{}\renewcommand{\headrulewidth}{0pt}
  \fancyfoot[L]{\raisebox{-0.32\height}{\includegraphics[height=0.78cm]{logo-gva.png}}\hspace{5mm}%
    \raisebox{-0.32\height}{\includegraphics[height=0.78cm]{logo-vermellar.png}}}
  \fancyfoot[R]{\scriptsize\color{muted}\href{https://creativecommons.org/licenses/by-nc-sa/4.0/deed.ca}{\ccbyncsa}\ \ CC BY-NC-SA 4.0 · fentmates.org\qquad\small\thepage}}
% la resta: el símbol a l'esquerra i el curs a la dreta, amb el títol al centre com abans
\fancyhead[L]{\raisebox{-0.3\height}{\includegraphics[height=0.62cm]{\fm@marca}}}
\fancyhead[R]{\small\bfseries\color{cursfosc}\fm@cursnom}
\renewcommand{\headrule}{\color{curs!45}\hrule height \headrulewidth}
\makeatother

\fmcurs{1eso}

\begin{document}
\thispagestyle{primera}

\capcalera{Matemáticas}{1\textsuperscript{o} ESO}{Ficha de ejercicios · Unidad 4}{Fracciones}
\vspace{10pt}
Nombre y apellidos: \rule{6.6cm}{0.4pt}\hfill Grupo: \rule{1.6cm}{0.4pt}\hfill Fecha: \rule{2.2cm}{0.4pt}
\vspace{8pt}

\begin{tcolorbox}[colback=soft,colframe=soft,arc=2mm,left=3mm,right=3mm,top=1.5mm,bottom=1.5mm]
\small Ejercicios graduados por niveles de dificultad:\quad \estrelles{1}\ básico \qquad \estrelles{2}\ medio \qquad \estrelles{3}\ avanzado.\\[2pt] Resuélvelos en el cuaderno mostrando todos los pasos y da los resultados en forma de fracción irreducible. Las figuras y las rectas se pueden completar directamente en la ficha. Haz los cálculos sin calculadora.
\end{tcolorbox}

% =====================================================
\seccio{1}{El concepto de fracción}

\ex{1} ¿Qué fracción representa la parte pintada de cada figura?\begin{tasks}(3)
  \task \begin{tikzpicture}[scale=0.42,baseline=(current bounding box.center)]\fill[accent!75] (0,1) rectangle ++(1,1);\fill[accent!75] (1,1) rectangle ++(1,1);\fill[accent!75] (2,1) rectangle ++(1,1);\fill[accent!75] (3,1) rectangle ++(1,1);\draw[thick] (0,0) grid (5,2);\end{tikzpicture}
  \task \begin{tikzpicture}[scale=0.42,baseline=(current bounding box.center)]\fill[accent!75] (0,1) rectangle ++(1,1);\fill[accent!75] (1,1) rectangle ++(1,1);\fill[accent!75] (0,0) rectangle ++(1,1);\draw[thick] (0,0) grid (2,2);\end{tikzpicture}
  \task \begin{tikzpicture}[baseline=(current bounding box.center)]\fill[accent!75] (0,0) -- (90.00:0.75) arc[start angle=90.00,end angle=38.57,radius=0.75] -- cycle;\draw[thick] (0,0) circle (0.75);\draw[thick] (0,0) -- (90.00:0.75);\draw[thick] (0,0) -- (38.57:0.75);\draw[thick] (0,0) -- (-12.86:0.75);\draw[thick] (0,0) -- (-64.29:0.75);\draw[thick] (0,0) -- (-115.71:0.75);\draw[thick] (0,0) -- (-167.14:0.75);\draw[thick] (0,0) -- (-218.57:0.75);\end{tikzpicture}
  \task \begin{tikzpicture}[scale=0.42,baseline=(current bounding box.center)]\fill[accent!75] (0,1) rectangle ++(1,1);\fill[accent!75] (1,1) rectangle ++(1,1);\draw[thick] (0,0) grid (2,2);\end{tikzpicture}
  \task \begin{tikzpicture}[baseline=(current bounding box.center)]\fill[accent!75] (0,0) -- (90.00:0.75) arc[start angle=90.00,end angle=38.57,radius=0.75] -- cycle;\fill[accent!75] (0,0) -- (38.57:0.75) arc[start angle=38.57,end angle=-12.86,radius=0.75] -- cycle;\fill[accent!75] (0,0) -- (-12.86:0.75) arc[start angle=-12.86,end angle=-64.29,radius=0.75] -- cycle;\fill[accent!75] (0,0) -- (-64.29:0.75) arc[start angle=-64.29,end angle=-115.71,radius=0.75] -- cycle;\fill[accent!75] (0,0) -- (-115.71:0.75) arc[start angle=-115.71,end angle=-167.14,radius=0.75] -- cycle;\draw[thick] (0,0) circle (0.75);\draw[thick] (0,0) -- (90.00:0.75);\draw[thick] (0,0) -- (38.57:0.75);\draw[thick] (0,0) -- (-12.86:0.75);\draw[thick] (0,0) -- (-64.29:0.75);\draw[thick] (0,0) -- (-115.71:0.75);\draw[thick] (0,0) -- (-167.14:0.75);\draw[thick] (0,0) -- (-218.57:0.75);\end{tikzpicture}
  \task \begin{tikzpicture}[scale=0.42,baseline=(current bounding box.center)]\fill[accent!75] (0,1) rectangle ++(1,1);\fill[accent!75] (1,1) rectangle ++(1,1);\fill[accent!75] (2,1) rectangle ++(1,1);\fill[accent!75] (0,0) rectangle ++(1,1);\draw[thick] (0,0) grid (3,2);\end{tikzpicture}
\end{tasks}\recorda{El \textbf{denominador} indica en cuántas partes \textbf{iguales} dividimos la unidad; el \textbf{numerador}, cuántas tomamos.}

\ex{1} Indica el numerador y el denominador de cada fracción y escribe cómo se lee:\begin{tasks}(4)
  \task $\dfrac{3}{5}$
  \task $\dfrac{4}{6}$
  \task $\dfrac{12}{15}$
  \task $\dfrac{4}{3}$
\end{tasks}

\ex{1} Pinta en cada figura la fracción indicada:\begin{tasks}(3)
  \task $\dfrac{3}{5}$\ \ \begin{tikzpicture}[baseline=(current bounding box.center)]\draw (0.700,0)--(0.700,0.55);\draw (1.400,0)--(1.400,0.55);\draw (2.100,0)--(2.100,0.55);\draw (2.800,0)--(2.800,0.55);\draw[thick] (0,0) rectangle (3.5,0.55);\end{tikzpicture}
  \task $\dfrac{5}{8}$\ \ \begin{tikzpicture}[scale=0.42,baseline=(current bounding box.center)]\draw[thick] (0,0) grid (4,2);\end{tikzpicture}
  \task $\dfrac{2}{3}$\ \ \begin{tikzpicture}[baseline=(current bounding box.center)]\draw[thick] (0,0) circle (0.65);\draw[thick] (0,0) -- (90.00:0.65);\draw[thick] (0,0) -- (-30.00:0.65);\draw[thick] (0,0) -- (-150.00:0.65);\end{tikzpicture}
\end{tasks}

\ex{1} Representa en la recta numérica las fracciones $\dfrac{6}{10}$ y $\dfrac{7}{4}$.\par\smallskip{\centering \begin{tikzpicture}[baseline=0]\draw[-latex,thick] (-0.3,0)--(9.50,0);\draw[thick] (0.00,0.18)--(0.00,-0.18) node[below=1pt,font=\small]{$0$};\draw[thick] (4.50,0.18)--(4.50,-0.18) node[below=1pt,font=\small]{$1$};\draw[thick] (9.00,0.18)--(9.00,-0.18) node[below=1pt,font=\small]{$2$};\end{tikzpicture}\par}

\ex{2} Clasifica las fracciones en propias, impropias o iguales a la unidad. Escribe las impropias como unidades enteras más una fracción propia.\begin{tasks}(6)
  \task $\dfrac{7}{3}$
  \task $\dfrac{5}{8}$
  \task $\dfrac{9}{9}$
  \task $\dfrac{13}{4}$
  \task $\dfrac{2}{11}$
  \task $\dfrac{20}{6}$
\end{tasks}

\ex{2} ¿Qué fracción representa?\begin{tasks}(2)
  \task 15 minutos de una hora
  \task 4 meses de un año
  \task 5 días de una semana
  \task 300 gramos de 1 kg
\end{tasks}

\ex{3} Escribe todas las fracciones de denominador $9$ que son impropias pero menores que $2$. ¿Qué fracciones de denominador $6$ hay entre $2$ y $3$ (sin contarlos)?


% =====================================================
\seccio{2}{Fracción de un número}

\ex{1} Calcula:\begin{tasks}(3)
  \task $\dfrac{5}{6}$ de $48$
  \task $\dfrac{4}{7}$ de $490$
  \task $\dfrac{3}{5}$ de $200$
  \task $\dfrac{7}{8}$ de $32$
  \task $\dfrac{2}{5}$ de $60$
  \task $\dfrac{3}{6}$ de $180$
\end{tasks}\recorda{$\dfrac{3}{5}$ de $40$: dividimos entre el denominador y multiplicamos por el numerador: \ $(40:5)\cdot3=24$.}

\ex{1} En una clase de 30 alumnos, $\dfrac{2}{5}$ son chicos. ¿Cuántos chicos y cuántas chicas hay?

\ex{2} Un depósito de $2400$ litros está lleno hasta los $\dfrac{3}{8}$. ¿Cuántos litros tiene? ¿Cuántos faltan para llenarlo?

\ex{2} Calcula el número:\begin{tasks}(2)
  \task Los $\dfrac{3}{4}$ de un número son $36$.
  \task Los $\dfrac{2}{7}$ de un número son $18$.
\end{tasks}

\ex{2} Pablo se ha gastado $\dfrac{3}{5}$ de la paga y todavía le quedan $14$\,€. ¿Cuánto dinero tenía?

\ex{3} Calcula $\dfrac{2}{3}$ de $\dfrac{3}{4}$ de $360$.

\ex{3} Van de excursión 84 alumnos: $\dfrac{3}{7}$ son de 1\textsuperscript{o}, $\dfrac{1}{4}$ de 2\textsuperscript{o} y el resto de 3\textsuperscript{o}. ¿Cuántos hay de cada curso?


% =====================================================
\seccio{3}{Porcentajes y fracciones}

\ex{1} Calcula:\begin{tasks}(3)
  \task $20\,\%$ de $180$
  \task $45\,\%$ de $80$
  \task $12\,\%$ de $350$
  \task $85\,\%$ de $1300$
  \task $50\,\%$ de $64$
  \task $25\,\%$ de $120$
\end{tasks}\recorda{$p\,\%=\dfrac{p}{100}$. \ El $p\,\%$ de $N$ es $\dfrac{p\cdot N}{100}$.}

\ex{1} Escribe cada porcentaje como fracción irreducible:\begin{tasks}(4)
  \task $50\,\%$
  \task $25\,\%$
  \task $75\,\%$
  \task $20\,\%$
  \task $10\,\%$
  \task $40\,\%$
  \task $5\,\%$
  \task $35\,\%$
\end{tasks}

\ex{1} Expresa en porcentaje:\begin{tasks}(6)
  \task $\dfrac{1}{2}$
  \task $\dfrac{3}{4}$
  \task $\dfrac{2}{5}$
  \task $\dfrac{7}{10}$
  \task $\dfrac{9}{20}$
  \task $\dfrac{3}{25}$
\end{tasks}

\ex{2} Una chaqueta cuesta $80$\,€ y la rebajan un $25\,\%$. ¿Cuántos euros nos descuentan? ¿Cuánto pagamos?

\ex{2} En un instituto de 650 alumnos, el $52\,\%$ son chicas. ¿Cuántas chicas y cuántos chicos hay?

\ex{3} En una provincia de $378\,000$ personas, el $7\,\%$ son de Marruecos, el $25\,\%$ del este de Europa, el $12\,\%$ de Alemania y el resto son de aquí. ¿Qué porcentaje son de aquí? ¿Cuántas personas hay de cada grupo?

\ex{3} En una clase de 25 alumnos, 8 han suspendido el examen. ¿Qué porcentaje ha aprobado?


% =====================================================
\seccio{4}{Fracciones equivalentes}

\ex{1} Escribe, por amplificación, cuatro fracciones equivalentes a cada una:\begin{tasks}(6)
  \task $\dfrac{1}{4}$
  \task $\dfrac{2}{5}$
  \task $\dfrac{5}{15}$
  \task $\dfrac{7}{13}$
  \task $\dfrac{11}{17}$
  \task $\dfrac{1}{10}$
\end{tasks}\recorda{\textbf{Amplificar}: multiplicar numerador y denominador por el mismo número. \textbf{Simplificar}: dividirlos entre un divisor común. La fracción \textbf{irreducible} se obtiene dividiendo entre el m.c.d.}

\ex{1} Halla dos fracciones simplificadas de cada una:\begin{tasks}(5)
  \task $\dfrac{4}{12}$
  \task $\dfrac{12}{20}$
  \task $\dfrac{66}{84}$
  \task $\dfrac{15}{45}$
  \task $\dfrac{50}{100}$
\end{tasks}

\ex{1} ¿Son irreducibles?\begin{tasks}(3)
  \task $\dfrac{25}{35}$
  \task $\dfrac{5}{7}$
  \task $\dfrac{98}{49}$
\end{tasks}

\ex{2} Simplifica hasta llegar a la fracción irreducible:\begin{tasks}(5)
  \task $\dfrac{4}{6}$
  \task $\dfrac{4}{20}$
  \task $\dfrac{3}{12}$
  \task $\dfrac{2}{6}$
  \task $\dfrac{35}{40}$
  \task $\dfrac{75}{100}$
  \task $\dfrac{4}{16}$
  \task $\dfrac{14}{21}$
  \task $\dfrac{10}{15}$
  \task $\dfrac{25}{45}$
  \task $\dfrac{13}{52}$
  \task $\dfrac{33}{34}$
  \task $\dfrac{705}{120}$
  \task $\dfrac{170}{105}$
  \task $\dfrac{256}{512}$
\end{tasks}

\ex{2} Calcula la fracción irreducible de cada una e indica cuáles son equivalentes entre sí:\begin{tasks}(3)
  \task $\dfrac{54}{24}$
  \task $\dfrac{40}{88}$
  \task $\dfrac{81}{36}$
  \task $\dfrac{72}{32}$
  \task $\dfrac{75}{185}$
  \task $\dfrac{50}{110}$
\end{tasks}

\ex{3} Halla una fracción equivalente a $\dfrac{15}{24}$ que tenga denominador $16$, y una equivalente a $\dfrac{5}{8}$ que tenga numerador $80$.

\ex{3} Escribe todas las fracciones equivalentes a $\dfrac{2}{3}$ que tienen el denominador entre $20$ y $30$.


% =====================================================
\seccio{5}{Comprobar si dos fracciones son equivalentes}

\ex{1} Comprueba si son equivalentes:\begin{tasks}(3)
  \task $\dfrac{8}{32}$ y $\dfrac{3}{12}$
  \task $\dfrac{6}{15}$ y $\dfrac{26}{34}$
  \task $\dfrac{7}{23}$ y $\dfrac{14}{46}$
  \task $\dfrac{1}{2}$ y $\dfrac{15}{30}$
  \task $\dfrac{68}{52}$ y $\dfrac{17}{13}$
  \task $\dfrac{3}{29}$ y $\dfrac{33}{67}$
\end{tasks}\recorda{$\dfrac{a}{b}=\dfrac{c}{d}$ \ si y solo si \ $a\cdot d=b\cdot c$ \ (productos cruzados).}

\ex{1} Calcula el valor de $x$ para que las fracciones sean equivalentes:\begin{tasks}(5)
  \task $\dfrac{3}{14}=\dfrac{x}{42}$
  \task $\dfrac{6}{8}=\dfrac{24}{x}$
  \task $\dfrac{x}{5}=\dfrac{12}{20}$
  \task $\dfrac{14}{3}=\dfrac{x}{21}$
  \task $\dfrac{5}{11}=\dfrac{15}{x}$
\end{tasks}

\ex{2} Completa para que sean equivalentes: \quad $\dfrac{2}{5}=\dfrac{\square}{20}=\dfrac{12}{\square}=\dfrac{\square}{45}$

\ex{2} ¿Cuál de estas fracciones no es equivalente a las demás? \quad $\dfrac{6}{9},\ \dfrac{10}{15},\ \dfrac{14}{20},\ \dfrac{16}{24}$

\ex{3} Halla $x$ e $y$ para que \ $\dfrac{x}{12}=\dfrac{15}{20}=\dfrac{9}{y}$.

\ex{3} Halla una fracción equivalente a $\dfrac{3}{5}$ que tenga el denominador $24$ unidades mayor que el numerador.


% =====================================================
\seccio{6}{Reducción a común denominador}

\ex{1} Reduce a común denominador:\begin{tasks}(4)
  \task $\dfrac{1}{5},\ \dfrac{3}{7}$
  \task $\dfrac{2}{9},\ \dfrac{1}{6}$
  \task $\dfrac{3}{4},\ \dfrac{3}{8}$
  \task $\dfrac{2}{7},\ \dfrac{1}{4}$
  \task $\dfrac{5}{21},\ \dfrac{3}{2}$
  \task $\dfrac{3}{5},\ \dfrac{1}{10}$
  \task $\dfrac{1}{4},\ \dfrac{1}{6}$
  \task $\dfrac{3}{7},\ 5$
\end{tasks}\recorda{Nuevo denominador: el \textbf{m.c.m.} de los denominadores. Para cada fracción, dividimos el m.c.m. entre su denominador y multiplicamos el numerador por el resultado.}

\ex{2} Reduce a común denominador:\begin{tasks}(4)
  \task $\dfrac{7}{24},\ \dfrac{3}{16}$
  \task $\dfrac{1}{2},\ \dfrac{2}{3},\ \dfrac{5}{8}$
  \task $\dfrac{7}{2},\ \dfrac{6}{15},\ \dfrac{3}{10}$
  \task $\dfrac{5}{8},\ \dfrac{1}{2},\ \dfrac{7}{12}$
  \task $\dfrac{4}{7},\ \dfrac{1}{3},\ \dfrac{5}{14}$
  \task $\dfrac{3}{18},\ \dfrac{4}{30},\ \dfrac{7}{36}$
  \task $\dfrac{6}{7},\ \dfrac{2}{14},\ \dfrac{12}{28}$
  \task $\dfrac{5}{3},\ \dfrac{5}{12},\ \dfrac{3}{4},\ \dfrac{9}{20}$
\end{tasks}

\ex{2} Escribe las fracciones $\dfrac{5}{6},\ \dfrac{7}{12},\ \dfrac{4}{9}$ con denominador $36$.

\ex{3} Simplifica primero y después reduce a común denominador con el denominador más pequeño posible: \ $\dfrac{6}{9},\ \dfrac{10}{15},\ \dfrac{9}{12}$.


% =====================================================
\seccio{7}{Comparar y ordenar fracciones}

\ex{1} Escribe una fracción mayor y una menor que la dada:\begin{tasks}(3)
  \task $\dfrac{5}{7}$
  \task $\dfrac{4}{9}$
  \task $\dfrac{3}{13}$
\end{tasks}

\ex{1} Compara (escribe $<$, $>$ o $=$):\begin{tasks}(4)
  \task $\dfrac{9}{5}\ \square\ \dfrac{9}{8}$
  \task $\dfrac{6}{11}\ \square\ \dfrac{10}{11}$
  \task $\dfrac{3}{7}\ \square\ \dfrac{3}{4}$
  \task $\dfrac{4}{8}\ \square\ \dfrac{2}{16}$
  \task $\dfrac{2}{5}\ \square\ \dfrac{5}{6}$
  \task $\dfrac{1}{3}\ \square\ \dfrac{3}{7}$
  \task $\dfrac{8}{9}\ \square\ \dfrac{11}{12}$
  \task $\dfrac{3}{4}\ \square\ \dfrac{24}{32}$
\end{tasks}\recorda{Mismo denominador: es mayor la de numerador mayor. Mismo numerador: es mayor la de denominador menor. Si no, reduce a común denominador.}

\ex{2} Ordena de menor a mayor:\begin{tasks}(3)
  \task $\dfrac{3}{5},\ \dfrac{2}{7},\ \dfrac{11}{2}$
  \task $\dfrac{4}{9},\ \dfrac{5}{6},\ \dfrac{2}{3}$
  \task $\dfrac{5}{4},\ \dfrac{2}{6},\ \dfrac{1}{10}$
  \task $\dfrac{18}{3},\ \dfrac{14}{7},\ \dfrac{23}{21}$
  \task $\dfrac{2}{12},\ \dfrac{4}{15},\ \dfrac{3}{24}$
  \task $\dfrac{12}{14},\ \dfrac{6}{7},\ \dfrac{3}{28}$
\end{tasks}

\ex{2} Ordena de menor a mayor:\begin{tasks}(3)
  \task $\dfrac{4}{5},\ \dfrac{2}{8},\ \dfrac{3}{10},\ \dfrac{7}{20}$
  \task $\dfrac{5}{2},\ \dfrac{2}{3},\ \dfrac{4}{7},\ \dfrac{5}{4}$
  \task $\dfrac{7}{12},\ \dfrac{2}{6},\ \dfrac{1}{4},\ \dfrac{3}{2}$
\end{tasks}

\ex{2} Marca en la recta las fracciones $\dfrac{1}{4},\ \dfrac{3}{4},\ \dfrac{5}{4},\ \dfrac{3}{2}$.\par\smallskip{\centering \begin{tikzpicture}[baseline=0]\draw[-latex,thick] (-0.3,0)--(9.50,0);\draw[muted] (1.125,0.1)--(1.125,-0.1);\draw[muted] (2.250,0.1)--(2.250,-0.1);\draw[muted] (3.375,0.1)--(3.375,-0.1);\draw[muted] (5.625,0.1)--(5.625,-0.1);\draw[muted] (6.750,0.1)--(6.750,-0.1);\draw[muted] (7.875,0.1)--(7.875,-0.1);\draw[thick] (0.00,0.18)--(0.00,-0.18) node[below=1pt,font=\small]{$0$};\draw[thick] (4.50,0.18)--(4.50,-0.18) node[below=1pt,font=\small]{$1$};\draw[thick] (9.00,0.18)--(9.00,-0.18) node[below=1pt,font=\small]{$2$};\end{tikzpicture}\par}

\ex{3} Escribe una fracción comprendida entre $\dfrac{2}{5}$ y $\dfrac{3}{5}$, y otra entre $\dfrac{1}{3}$ y $\dfrac{1}{2}$.

\ex{3} Tres amigos leen el mismo libro: Ana ha leído $\dfrac{3}{4}$, Gabriel $\dfrac{5}{8}$ y Clara $\dfrac{7}{10}$. ¿Quién ha leído más? Ordénalos.


% =====================================================
\seccio{8}{Sumas y restas}

\ex{1} Calcula y simplifica el resultado:\begin{tasks}(3)
  \task $\dfrac{4}{7}+\dfrac{8}{7}$
  \task $\dfrac{2}{3}+\dfrac{1}{5}$
  \task $\dfrac{5}{2}-\dfrac{1}{8}$
  \task $\dfrac{3}{6}-\dfrac{4}{12}$
  \task $\dfrac{5}{8}+\dfrac{1}{5}$
  \task $\dfrac{14}{15}-\dfrac{2}{3}$
\end{tasks}\recorda{Si los denominadores son diferentes, primero reducimos a \textbf{común denominador}. \textbf{Nunca} sumamos los denominadores: $\dfrac{1}{2}+\dfrac{1}{3}=\dfrac{3}{6}+\dfrac{2}{6}=\dfrac{5}{6}$.}

\ex{1} Calcula y simplifica (los números naturales tienen denominador 1):\begin{tasks}(3)
  \task $\dfrac{3}{4}+\dfrac{6}{8}$
  \task $4+\dfrac{25}{9}$
  \task $\dfrac{25}{8}-2$
  \task $\dfrac{3}{5}+1$
  \task $5-\dfrac{3}{25}$
  \task $\dfrac{13}{6}+\dfrac{2}{7}$
\end{tasks}

\ex{2} Calcula y simplifica el resultado:\begin{tasks}(3)
  \task $\dfrac{12}{15}+\dfrac{8}{9}$
  \task $\dfrac{7}{12}-\dfrac{3}{10}$
  \task $\dfrac{7}{6}+\dfrac{8}{12}$
  \task $\dfrac{18}{5}-\dfrac{3}{10}$
  \task $\dfrac{5}{3}-\dfrac{9}{24}$
  \task $\dfrac{35}{36}+\dfrac{2}{9}$
\end{tasks}

\ex{2} Calcula y simplifica:\begin{tasks}(2)
  \task $\dfrac{4}{5}+\dfrac{2}{7}-\dfrac{1}{2}$
  \task $\dfrac{6}{8}+\dfrac{4}{3}-\dfrac{8}{12}$
  \task $\dfrac{3}{2}+\dfrac{6}{7}-\dfrac{2}{4}$
  \task $6-\dfrac{1}{2}-\dfrac{4}{12}$
  \task $\dfrac{2}{3}+\dfrac{1}{5}-\dfrac{1}{6}$
  \task $\dfrac{4}{3}+\dfrac{1}{6}-\dfrac{2}{7}$
\end{tasks}

\ex{3} Calcula y simplifica:\begin{tasks}(2)
  \task $\dfrac{2}{10}+\dfrac{5}{60}-\dfrac{1}{30}$
  \task $\dfrac{12}{5}+2-\dfrac{1}{25}$
  \task $\dfrac{6}{9}+\dfrac{1}{3}-\dfrac{2}{6}$
  \task $\dfrac{10}{3}-\dfrac{2}{6}-\dfrac{4}{5}+\dfrac{1}{2}$
  \task $4-\dfrac{1}{2}+\dfrac{1}{3}+2$
  \task $\dfrac{7}{9}+\dfrac{1}{3}-\dfrac{2}{5}$
\end{tasks}

\ex{3} Llenamos una piscina en tres días: el primer día, $\dfrac{1}{4}$; el segundo, $\dfrac{2}{5}$, y el tercero, $\dfrac{1}{6}$. ¿Qué fracción de la piscina falta por llenar?


% =====================================================
\seccio{9}{Multiplicación y división}

\ex{1} Calcula y simplifica:\begin{tasks}(3)
  \task $\dfrac{6}{4}\cdot \dfrac{2}{5}$
  \task $\dfrac{3}{2}\cdot \dfrac{2}{3}$
  \task $2\cdot \dfrac{5}{6}$
  \task $\dfrac{3}{2}\cdot \dfrac{4}{7}\cdot 3$
  \task $\dfrac{3}{4}\cdot \dfrac{10}{2}\cdot \dfrac{1}{6}$
  \task $\dfrac{4}{3}\cdot \dfrac{2}{8}\cdot \dfrac{6}{2}$
\end{tasks}\recorda{$\dfrac{a}{b}\cdot\dfrac{c}{d}=\dfrac{a\cdot c}{b\cdot d}$. \qquad $\dfrac{a}{b}:\dfrac{c}{d}=\dfrac{a\cdot d}{b\cdot c}$ \ (productos cruzados). \ La \textbf{inversa} de $\dfrac{a}{b}$ es $\dfrac{b}{a}$.}

\ex{1} Calcula (la fracción de una fracción es un producto):\begin{tasks}(4)
  \task $\dfrac{2}{3}$ de $\dfrac{3}{5}$
  \task $\dfrac{4}{10}$ de $\dfrac{1}{2}$
  \task $\dfrac{4}{3}$ de $\dfrac{6}{2}$
  \task $\dfrac{3}{5}$ de $\dfrac{5}{7}$
\end{tasks}

\ex{1} ¿Qué parejas son fracciones inversas? (Su producto tiene que ser 1.)\begin{tasks}(3)
  \task $\dfrac{2}{5},\ \dfrac{5}{3}$
  \task $\dfrac{4}{3},\ \dfrac{12}{16}$
  \task $\dfrac{7}{9},\ \dfrac{9}{7}$
  \task $\dfrac{2}{4},\ \dfrac{4}{2}$
  \task $\dfrac{5}{6},\ \dfrac{36}{30}$
  \task $\dfrac{12}{13},\ \dfrac{26}{12}$
\end{tasks}

\ex{1} Escribe la inversa de:\begin{tasks}(4)
  \task $\dfrac{4}{13}$
  \task $\dfrac{3}{7}$
  \task $\dfrac{15}{8}$
  \task $6$
\end{tasks}

\ex{2} Calcula y simplifica (con divisiones seguidas, de izquierda a derecha):\begin{tasks}(3)
  \task $\dfrac{4}{3}:\dfrac{5}{6}$
  \task $\dfrac{5}{4}:2$
  \task $\dfrac{4}{8}:\dfrac{6}{3}$
  \task $\dfrac{2}{6}:\dfrac{5}{4}:\dfrac{1}{3}$
  \task $\dfrac{3}{4}:\dfrac{12}{5}:\dfrac{2}{5}$
  \task $6:\dfrac{2}{3}:\dfrac{1}{5}$
\end{tasks}

\ex{2} Calcula y simplifica:\begin{tasks}(2)
  \task $\dfrac{3}{4}\cdot \dfrac{5}{8}:\dfrac{2}{5}$
  \task $\dfrac{5}{4}:\dfrac{1}{3}\cdot \dfrac{6}{5}$
  \task $\dfrac{5}{12}\cdot \dfrac{3}{2}:\dfrac{1}{6}$
  \task $\dfrac{6}{5}:\dfrac{7}{4}\cdot \dfrac{2}{10}$
\end{tasks}

\ex{3} Con $12$ litros de zumo, ¿cuántas botellas de $\dfrac{3}{4}$ de litro podemos llenar? ¿Y qué fracción hay que multiplicar por $\dfrac{3}{4}$ para obtener $\dfrac{9}{10}$?


% =====================================================
\seccio{10}{Operaciones combinadas}

\ex{2} Calcula y simplifica:\begin{tasks}(2)
  \task $\dfrac{5}{2}+\dfrac{1}{3}\cdot \dfrac{4}{5}$
  \task $\dfrac{21}{4}-\dfrac{2}{5}:\dfrac{4}{3}$
  \task $\dfrac{4}{3}:\dfrac{2}{6}+\dfrac{1}{3}\cdot \dfrac{5}{4}$
  \task $\dfrac{4}{3}\cdot \left(\dfrac{2}{5}+\dfrac{1}{3}\right)$
  \task $\left(\dfrac{2}{5}+1\right):\dfrac{8}{3}$
  \task $\left(\dfrac{5}{4}+1\right)\cdot \dfrac{6}{7}$
\end{tasks}\recorda{Orden: primero, \textbf{paréntesis}; después, \textbf{multiplicaciones y divisiones} (de izquierda a derecha); por último, \textbf{sumas y restas} (de izquierda a derecha).}

\ex{2} Calcula y simplifica:\begin{tasks}(2)
  \task $\left(\dfrac{2}{5}+\dfrac{3}{7}\right)\cdot \dfrac{1}{2}-\dfrac{4}{10}$
  \task $\left(\dfrac{5}{4}+\dfrac{2}{3}+\dfrac{1}{12}\right)\cdot \dfrac{5}{4}-1$
  \task $\dfrac{5}{4}:\left(\dfrac{6}{4}-1\right)+\dfrac{2}{5}$
  \task $\dfrac{4}{3}-\left(\dfrac{1}{2}+\dfrac{2}{7}\right)+\dfrac{4}{5}\cdot \dfrac{2}{4}$
\end{tasks}

\ex{3} Calcula y simplifica:\begin{tasks}(2)
  \task $\left(\dfrac{1}{4}+\dfrac{2}{3}-\dfrac{1}{12}\right)\cdot 3+\left(\dfrac{3}{2}-\dfrac{5}{4}\right):\dfrac{1}{2}$
  \task $4\cdot \left(\dfrac{2}{3}+\dfrac{4}{6}-1\right)+\dfrac{2}{3}$
  \task $\left(\dfrac{3}{2}+\dfrac{1}{7}\right)\cdot \left(\dfrac{5}{3}-\dfrac{7}{15}\right)$
  \task $\dfrac{12}{10}:\dfrac{5}{4}+\dfrac{1}{2}\cdot \dfrac{13}{5}$
  \task $\left(\dfrac{4}{3}+\dfrac{2}{5}\cdot \dfrac{1}{3}\right):\dfrac{4}{12}$
  \task $\dfrac{3}{2}:\dfrac{1}{4}\cdot \dfrac{4}{20}+\dfrac{3}{2}\cdot \dfrac{7}{5}$
\end{tasks}

\ex{3} Calcula y simplifica:\begin{tasks}(1)
  \task $\left(\dfrac{5}{4}+\dfrac{2}{3}-1\right)\cdot \dfrac{3}{5}+\left(\dfrac{2}{5}+\dfrac{4}{3}\right)\cdot 2$
  \task $\dfrac{3}{8}+\dfrac{4}{5}\cdot \dfrac{1}{3}+\left(\dfrac{2}{6}-\dfrac{1}{4}\right)\cdot 3$
\end{tasks}


% =====================================================
\seccio{11}{Problemas con fracciones}

\ex{1} De una cena, Juan se ha comido $\dfrac{1}{7}$, Elena $\dfrac{2}{5}$ y Maite $\dfrac{1}{8}$.\begin{tasks}(1)
  \task ¿Cuánto se han comido entre los tres?
  \task ¿Cuánto ha sobrado?
  \task ¿Quién ha comido más y quién menos?
\end{tasks}\recorda{Lee el enunciado dos veces, haz un \textbf{dibujo} (una barra que representa el total) y, al final, \textbf{comprueba} que el resultado tiene sentido.}

\ex{1} De un pastel de 24 porciones nos hemos comido $\dfrac{5}{8}$. ¿Cuántas porciones quedan?

\ex{2} De una pizza, Carlos se come $\dfrac{3}{6}$, Miguel $\dfrac{3}{21}$ y Lorena $\dfrac{4}{14}$.\begin{tasks}(1)
  \task ¿Qué fracción de la pizza ha sobrado?
  \task Llega Marta con otra pizza igual. Como los demás ya están llenos, ella se come lo que había sobrado y $\dfrac{3}{5}$ de su pizza. ¿Cuánta pizza se ha comido ella?
\end{tasks}

\ex{2} En un pueblo de 5250 habitantes, la mitad tienen más de 65 años y $\dfrac{1}{5}$ tienen menos de 18.\begin{tasks}(1)
  \task ¿Qué fracción representan los que tienen entre 18 y 65 años?
  \task ¿Cuántos habitantes hay en cada franja de edad?
\end{tasks}

\ex{2} Un bidón de 60 litros está lleno. Primero se gastan $\dfrac{2}{4}$ y después $\dfrac{1}{3}$ de lo que quedaba. Finalmente, el bidón vuelca y se derrama la mitad de lo que hay. ¿Cuántos litros quedan al final?

\ex{3} Una familia gasta $\dfrac{2}{6}$ del sueldo en la hipoteca, $\dfrac{1}{5}$ en comida, $\dfrac{2}{7}$ en el coche y el resto en otros gastos. ¿Qué fracción representan los otros gastos?

\ex{3} Una familia ingresa $5400$\,€ al mes. Gasta $\dfrac{2}{6}$ en la hipoteca, $\dfrac{1}{5}$ en comida y $\dfrac{1}{9}$ en el coche.\begin{tasks}(1)
  \task ¿Cuánto dinero gasta en cada cosa y cuánto le queda para otros gastos?
  \task ¿Qué fracción del total representa el dinero que sobra?
\end{tasks}

\ex{3} Un depósito está lleno hasta los $\dfrac{3}{4}$. Gastamos $\dfrac{2}{3}$ de lo que hay y quedan 50 litros. ¿Qué capacidad tiene el depósito?

\vfill
\begin{center}\color{muted}\small\faCheckCircle[regular]\ \ ¡Ánimo! Si te cuesta algún ejercicio, repasa los apuntes interactivos de la unidad.\end{center}

\end{document}
